% ============================================================================
% SECTION 4 --- Controlling molecular trajectories
%   * objective regions and Chebyshev preferences, scale fixed by a statistic
%     of the SUPERVISION distribution and never by an outcome
%   * dynamic retargeting, as a compact corollary
%   * pathwise constraints: rule-closed support restriction vs cumulative
%     constraints needing an augmented state
%   * Pareto exploration
%
% PROVENANCE NOTE FOR LATER EDITORS.  The preference construction follows
% docs/AMENDMENT_INVERSIONGNN_2OBJ.md, "ADDENDUM II -- SUPERSEDES the B_lambda
% construction above" (2026-08-19).  Addendum I's top-10% region construction
% is WITHDRAWN and must not be reintroduced: it trained the value on observed
% true-oracle hits, and its own diagnostic showed the five preferences
% collapsing to three regions.  If you are editing this section from an older
% document, you are editing from the withdrawn design.
% No number appears in this section.  Compressed to the 9-page limit.
% ============================================================================
\section{Controlling molecular trajectories}
\label{sec:trajectories}

\paragraph{Objective regions.}
\label{sec:regions}
Campaigns state two kinds of goal and each becomes a terminal desirability. A \emph{requirement} is a region, a property above a threshold and a similarity floor to the source, so $g_z=\one[x\in\Bset_z]$ and the corresponding value is the finite-budget reachability of \cref{sec:hphi}, the thresholds serving as the goal's coordinates. A \emph{preference} is a direction and needs a positive, smooth desirability (\cref{sec:preferences}). A threshold invented to turn a direction into a region is a hyperparameter that will be tuned, and with several objectives simultaneously low such a region degenerates, distinct preferences collapsing onto the same molecules.
With no threshold and no positive-example requirement nothing here can be slid after seeing a result, and three objects stay separated, which is the paper's thesis restated for the multi-objective case, where $\widehat F_\psi$ says \emph{what is desirable}, $R_\theta$ \emph{how molecules can plausibly move}, $h_\phi$ \emph{what is desirable in the future, given how molecules move}.

\begin{corollary}[Exact continuation after a goal change]
\label{cor:retarget}
Let a trajectory reach $x_\tau$ with $b'$ edits remaining under goal $z$, and let the goal become $z'$. Applying \cref{thm:doob} from $x_\tau$ with horizon $b'$, controlling the remaining $b'$ edits with the $g_{z'}$-value re-solves the finite-horizon problem exactly from $x_\tau$, and the terminal law is the $g_{z'}$-tilt of the base law started at $x_\tau$ with budget $b'$, in general \emph{not} the law obtained by using $g_{z'}$ from the source lead.
\end{corollary}

Committed work therefore need not be discarded when the objective moves, and comparing continuation against restart is substantive rather than tautological. The transform applies because $x_\tau$ is a molecule (\cref{prop:closure}), so $g_{z'}$ is decidable on it, whereas an endpoint-conditioned generator has no analogue. It does not assert that continuing is \emph{better}.

\paragraph{Pathwise constraints and Pareto exploration.}
\label{sec:pathwise}
A requirement holding at every inspected candidate is a property of the path, and support restriction makes one class exact by construction. For a \emph{rule-closed} predicate $P$ (a protected scaffold, a forbidden substructure, a size or charge window, a similarity floor) we delete violating marks before normalizing, $\A_P(x)=\{a\in\A(x):P(\T(x,a))\}$, so the restricted chain is again a chain on $\X$, \cref{thm:doob} applies unchanged, $P$ holds at \emph{every} committed molecule by induction rather than by filtering at the end, and a candidate rejected by $P$ is rejected before a computational property oracle is called on it.
A \emph{cumulative} constraint is not a predicate on the current molecule and needs the state augmented with the running statistic, and we report which of the two any constraint is.
Multi-objective design is then a loop over goals rather than scalarization weights. Set $z$ or $\lambda$ to an under-covered direction of a nondominated archive, continue a controlled trajectory toward it under the frozen $R_\theta$, and fold every committed molecule into the archive, which is meaningful because every committed molecule is a candidate.
Branching is what a restart-based optimizer cannot do. By \cref{cor:retarget}, several preference-conditioned continuations launch from one saved prefix, each the exact finite-horizon problem for its own preference on the remaining budget, a \emph{Pareto fan}.
